% page 298 du fac-simile (image p-297 du scan)
% bloc 152.4 x 242.0 mm ; marge gauche 8.0 mm
\pgc{-12.700mm}{0.000mm}{
\l{290}
\l{}
\l{\sou{komparar}. (\sou{trans.}) I. (komparar\cel{1}\sou{kun}:)\cel{3}Apudigar du o plu-}
\l{\cel{3}ra kozi por determinar lia punti di simileso e di ne-}
\l{\cel{3}simileso reciproka. - II. (komparar ad:) . Apudigar}
\l{\cel{3}por igar relatar egalesale. - III. Apudigar to quon onu}
\l{\cel{3}aludas, ad ulo analoga, por igar plu komprenebla, plu}
\l{\cel{3}sentebla, lo dicata. - EFIS.}
\l{}
\l{\sou{komparativo}.\cel{2}(\sou{gram.}) Modifikuro di la adjektivo, qua indi-}
\l{\cel{3}kas ke la qualeso expresata da lu pozitive esas egardata}
\l{\cel{3}ye grado superiora. - DEFIS.}
\l{}
\l{\sou{kompaso}. Instrumento ek du branchi de ligno o de metalo,}
\l{\cel{3}asemblita ye un ek lia extremaji, tale ke li povas inter-}
\l{\cel{3}forigar od inter-apudigar, por mezurar anguli, trasar}
\l{\cel{3}cirkli dimensionale inter-distanta. - EFIRS.}
\l{}
\l{\sou{kompatar}. (\sou{trans., pri, pro}).\cel{2}Partoprenar la sufro, la}
\l{\cel{3}chagreno, da altru. - EFIS.}
\l{}
\l{\sou{kompendio}.\cel{2}Rezumuro ek la ensemblo di cienco, di doktrino.}
\l{\cel{3}- DEFIRS.}
\l{}
\l{\sou{kompensar}.(\cel{1}\sou{trans., per}).\cel{2}Neutrigar (efekto) per efekto}
\l{\cel{3}equivalanta, inversa-sinse. - DEFIRS.}
\l{}
\l{\sou{kompetenta}.\cel{2}A qua esas atribuenda la yuro decidar : l. kon-}
\l{\cel{3}seque de autoritato lego-grantita ; 2. konseque de par-}
\l{\cel{3}konoco, par-savo, de parposedo di la temo koncernata. -}
\l{\cel{3}DEFIRS.}
\l{}
\l{\sou{kompilar}. (\sou{trans., ek, de)}.\cel{2}Asemblar aden un kolektajo,}
\l{\cel{3}aden un korpo (texti pri un temo komuna, cherpita de}
\l{\cel{3}fonti diversa). - DEFIRS.}
\l{}
\l{\sou{komplemento}.\cel{2}Lo adjuntenda a kozo por igar lu kompleta.}
\l{\cel{3}- DEFIRS.}
\l{}
\l{\sou{kompleta}. A qua mankas nulo ek la elementi ye qui lu kon-}
\l{\cel{3}sistas normale. - DEFIS.}
\l{}
\l{\sou{kompletorio}. (\sou{liturgio katol.}) La lasta parto di la oficio,}
\l{\cel{3}qua dicesas o kantesas pos la vespri. -}
\l{}
\l{\sou{komplexa}. Qua unionas en su plura elementi diversa. - DEFIRS}
\l{}
\l{\sou{komplexiono}. Ensemblo ek la elementi ye qui konsistas la}
\l{\cel{3}naturo korpala di individuo. - EF.}
\l{}
\l{\sou{komplezar}. (\sou{netrans}.)\cel{2}I. Plezar (ad ulu) per adaptar su}
\l{\cel{3}a lua gusti, a lua sentimenti. - II. Inklinesar a per-}
\l{\cel{3}turbar su por plezar (ad ulu). - EFIS.}
\l{}
\l{\sou{komplico}. La persono qua helpas (ulu) kriminar, deliktar.}
\l{\cel{3}- DEFIS.}
}
